
\section{极坐标轴}
\label{sec:polar}

{
\def\pgfplotsmanualcurlibrary{polar}
\tikzset{external/figure name/.add={}{polar_}}

\begin{pgfplotslibrary}{polar}
    本库用于绘制极坐标轴以及依赖极坐标的绘图类型，坐标由角度（以度为单位，
    或可选地以弧度为单位）与半径表示。
\end{pgfplotslibrary}


\subsection{极坐标轴}

\begin{environment}{{polaraxis}}
    |polar| 库提供 |polaraxis| 环境。
    在该环境内，所有坐标均应以 $(\mathit{angle},\mathit{radius})$ 的极坐标形式给出，
    即 $x$ 坐标始终为角度，$y$ 坐标始终为半径：
\end{environment}

\begin{codeexample}[]
\begin{tikzpicture}
\begin{polaraxis}
    \addplot coordinates {
        (0,1)
        (90,1)
        (180,1)
        (270,1)
    };
\end{polaraxis}
\end{tikzpicture}
\end{codeexample}

\begin{codeexample}[]
\begin{tikzpicture}
\begin{polaraxis}
    \addplot+ [domain=0:3] (360*x,x); % (angle,radius)
\end{polaraxis}
\end{tikzpicture}
\end{codeexample}

\begin{codeexample}[]
\begin{tikzpicture}
\begin{polaraxis}
    \addplot+ [
        mark=none,
        domain=0:720,
        samples=600,
    ]
        {sin(2*x)*cos(2*x)};
    % equivalent to (x,{sin(..)cos(..)}), i.e.
    % the expression is the RADIUS
\end{polaraxis}
\end{tikzpicture}
\end{codeexample}

极坐标轴支持 \PGFPlots{} 的大部分用户接口，例如 |legend entries|、各种坐标轴描述、
|xtick|/|ytick| 等：
%
\begin{codeexample}[]
\begin{tikzpicture}
\begin{polaraxis}[
    xtick={0,90,180,270},
    title=A polar axis,
]
    \addplot coordinates {(0,1) (45,1)};
        \addlegendentry{First}

    \addplot coordinates {(180,0.5) (0,0)};
        \addlegendentry{Second}
    \end{polaraxis}
\end{tikzpicture}
\end{codeexample}
%
\noindent 此外，你可以使用所有受支持的输入坐标方式（如
\verbpdfref{\addplot coordinates}、\verbpdfref{\addplot table}、
\verbpdfref{\addplot expression}）。唯一的差别在于：极坐标轴将（前两个）输入坐标
解释为 $(\mathit{angle\,in\,degrees},\mathit{radius})$ 形式的极坐标。

也可以使用 \verbpdfref{\addplot3}；此时第三个坐标将被忽略（不过它仍可通过
|point meta=z| 用作颜色数据）。示例见下文第~\ref{sec:polar:cart} 节。


\subsection{使用弧度代替角度}

初始配置中，角度（即每个输入坐标的 $x$ 分量）以度为单位。\PGFPlots{} 也支持通过
|data cs=polarrad| 开关以弧度提供角度：
%
\begin{codeexample}[]
\begin{tikzpicture}
\begin{polaraxis}[title={Degrees and/or Radians}]
    \addplot coordinates {
        (0,1) (90,1) (180,1) (270,1)
    };
        \addlegendentry{Deg}

    \addplot+ [
        data cs=polarrad,
    ] coordinates {
        (0,1.5) (pi/2,1.5)
        (pi,1.5) (pi*3/2,1.5)
    };
        \addlegendentry{Rad}
\end{polaraxis}
\end{tikzpicture}
\end{codeexample}

|data cs| 键的完整说明见第~\pageref{key:data:cs} 页；它告知 \PGFPlots{} 输入数据的
坐标系。随后 \PGFPlots{} 会自动将每个坐标变换到所需的坐标系（在本例中即
|data cs=polar|）。


\subsection{与笛卡尔坐标混用}
\label{sec:polar:cart}

与上述过程类似，你也可以在极坐标轴内提供笛卡尔坐标：只需通过 |data cs=cart|
告知 \PGFPlots{} 自动将其变换为极坐标形式：
%
\begin{codeexample}[]
\begin{tikzpicture}
\begin{polaraxis}[title=Cartesian Input]
    \addplot+ [
        data cs=cart,
    ] coordinates {
        (1,0) (0,1) (-1,0) (0,-1)
    };
\end{polaraxis}
\end{tikzpicture}
\end{codeexample}
%
\noindent 关于 |data cs| 键的更多细节见第~\pageref{key:data:cs} 页。

这也使得可以使用那些基于笛卡尔坐标工作的更复杂可视化技术。下面的例子使用
\verbpdfref{\addplot3} 对函数 $f\colon \R^2 \to \R$ 采样，计算 |contour| 等高线，
并将结果显示在 |polaraxis| 中：
%
\pgfplotsexpensiveexample
\begin{codeexample}[]
\begin{tikzpicture}
\begin{polaraxis}
    \addplot3 [
        contour lua,
        domain=-3:3,
        data cs=cart,
    ]
        {exp(-x^2-y^2)};
\end{polaraxis}
\end{tikzpicture}
\end{codeexample}
%
\noindent 具体过程为：在 $x,y \in [-3,3]$ 上对 $z=\exp(-x^2-y^2)$ 采样，然后在
$(x,y,z)$ 上计算等高线，再将所得三元组 $(x,y,z)$ 变换为极坐标 $(\alpha,r,z)$
（$z$ 保持不变）。最后，$z$ 坐标作为 |point meta| 用于决定颜色。

注意，\verbpdfref{\addplot3} 允许处理三维输入类型，但结果始终是二维的
（在 |polaraxis| 中，点位置的放置会忽略 $z$ 坐标）。不过，$z$ 坐标仍可用于
决定点的颜色（通过 |point meta=z|）。


\subsection{极坐标描述}

极坐标轴与其他绘图类型一样支持坐标轴标签。\footnote{注意：在 1.13 之前的版本中，
极坐标轴标签支持存在缺陷。请升级到 \PGFPlots{} 1.13 并使用
\texttt{compat=1.13} 或更高版本，以使用极坐标轴描述。}
此时，|xticklabel cs| 指角度（$x$ 轴）的标签位置，|yticklabel cs| 指半径（$y$ 轴）
的标签位置。

角度的默认坐标轴标签放置在 $45$ 度处，
%
\begin{codeexample}[]
\begin{tikzpicture}
\begin{polaraxis}[
    xlabel=x axis,
    ylabel=y axis,
]
    \addplot coordinates {
        (0,1) (90,1) (180,1) (270,1)
    };
\end{polaraxis}
\end{tikzpicture}
\end{codeexample}
%
并配置为 |sloped|\index{sloped Polar labels} 标签，采用如下默认配置：
%
\begin{codeexample}[code only]
    every axis x label/.style={
        at={(xticklabel cs:0.125)},
        sloped={at position=45},
        anchor=near ticklabel,
        near ticklabel at=45,
    },
    every axis y label/.style={
        at={(yticklabel cs:0.5)},
        anchor=near ticklabel,
    },
\end{codeexample}

下面的例子也展示了倾斜的 $x$ 标签。\footnote{注意，倾斜的极坐标标签同样需要
\PGFPlots{} 1.13 及 \texttt{compat=1.13} 或更新版本。}
%
\begin{codeexample}[]
\begin{tikzpicture}
\begin{polaraxis}[
    xlabel=x axis,
    ylabel=y axis,
    xticklabel style={sloped like x axis},
]
    \addplot coordinates {
        (0,1) (90,1) (180,1) (270,1)
    };
\end{polaraxis}
\end{tikzpicture}
\end{codeexample}


\subsection{特殊极坐标绘图类型}

\begin{plottype}{polar comb}
    |polar comb| 绘图处理器由 \Tikz{} 提供；它从原点向指定位置绘制路径，并在
    位置上放置 |mark|（与 |comb| 绘图处理器类似）。由于路径始终从原点出发，
    它特别适合 |polaraxis|：
    %
\begin{codeexample}[]
\begin{tikzpicture}
\begin{polaraxis}
    \addplot+ [
        polar comb,
    ] coordinates {
        (300,1)
        (20,0.3)
        (40,0.5)
        (120,1)
        (200,0.4)
    };
\end{polaraxis}
\end{tikzpicture}
\end{codeexample}
\end{plottype}


\subsection{局部极坐标轴}

|polar| 库也支持局部坐标轴。若你提供 |xmin|/|xmax|，即可限制坐标轴使用的角度：
%
\begin{codeexample}[]
\begin{tikzpicture}
\begin{polaraxis}[
    xmin=45,
    xmax=360,
]
    \addplot coordinates {
        (0,1) (90,1) (180,1) (270,1)
    };
\end{polaraxis}
\end{tikzpicture}
\end{codeexample}

目前，第一个角度必须小于第二个角度。但你可以利用周期性，按如下方式得到扇形：
%
\message{Overfull hbox is OK}%
\begin{codeexample}[]
\begin{tikzpicture}
    \begin{polaraxis}[xmin=90,xmax=270]
        \addplot coordinates { (0,1) (90,1) (180,1) (270,1) };
    \end{polaraxis}
\end{tikzpicture}~%
\begin{tikzpicture}
    \begin{polaraxis}[xmin=270,xmax=420]
        \addplot coordinates { (0,1) (90,1) (180,1) (270,1) };
    \end{polaraxis}
\end{tikzpicture}
\end{codeexample}
%
\noindent 类似地，显式提供 |ymin| 可将显示范围从 $0$ 缩小：
%
\begin{codeexample}[]
\begin{tikzpicture}
\begin{polaraxis}[
    ymin=0.3,
]
    \addplot coordinates {
        (0,1) (90,1) (180,1) (270,1)
    };
\end{polaraxis}
\end{tikzpicture}
\end{codeexample}

\noindent 手动修改 |xmin| 与 |xmax| 还可以用来移动 $y$ 轴线（即带有 |ytick| 和
|yticklabels| 的那条线）：
%
\begin{codeexample}[]
\begin{tikzpicture}
\begin{polaraxis}[
    xmin=45,
    xmax=405,
]
    \addplot coordinates {
        (0,1) (90,1) (180,1) (270,1)
    };
\end{polaraxis}
\end{tikzpicture}
\end{codeexample}
}
